% Created 2018-03-27 Tue 20:31
% Intended LaTeX compiler: pdflatex
\documentclass[presentation,smaller]{beamer}
\RequirePackage{etex}
\RequirePackage[l2tabu,orthodox]{nag}            %% Warn about obsolete commands and packages
\RequirePackage{amsmath,amsfonts,amssymb,amsthm} %% Math
\RequirePackage{ifxetex,ifluatex}                %% Detect XeTeX and LuaTeX
\RequirePackage{fixltx2e}                        %% provides \textsubscript
\RequirePackage{xspace}
\RequirePackage{graphicx}
\RequirePackage{comment}
\RequirePackage{url}
\RequirePackage{relsize}
\RequirePackage{booktabs}
\RequirePackage{tabularx}
\RequirePackage[normalem]{ulem}
\RequirePackage[all]{xy}
\RequirePackage{etoolbox}

%%%
%%% Code Listings
%%%

\RequirePackage{listings}
\lstdefinelanguage{Sage}[]{Python}{morekeywords={True,False,sage,cdef,cpdef,ctypedef,self},sensitive=true}

\lstset{frame=none,
  showtabs=False,
  showspaces=False,
  showstringspaces=False,
  commentstyle={\color{gray}},
  keywordstyle={\color{mLightBrown}\textbf},
  stringstyle ={\color{mDarkBrown}},
  frame=single,
  basicstyle=\tt\scriptsize\relax,
  backgroundcolor=\color{gray!190!black},
  inputencoding=utf8,
  literate={…}{{\ldots}}1,
  belowskip=0.0em,
}

\makeatletter
\patchcmd{\@verbatim}
  {\verbatim@font}
  {\verbatim@font\scriptsize}
  {}{}
\makeatother

%%%
%%% Tikz
%%%

\RequirePackage{tikz,pgfplots}

\usetikzlibrary{calc}
\usetikzlibrary{arrows}
\usetikzlibrary{automata}
\usetikzlibrary{positioning}
\usetikzlibrary{decorations.pathmorphing}
\usetikzlibrary{backgrounds}
\usetikzlibrary{fit,}
\usetikzlibrary{shapes.symbols}
\usetikzlibrary{chains}
\usetikzlibrary{shapes.geometric}
\usetikzlibrary{shapes.arrows}
\usetikzlibrary{graphs}

%% Cache

\ifdefined\tikzcaching  % chktex 1
  \usetikzlibrary{external}
  \tikzexternalize[prefix=build/]
  \tikzset{external/up to date check=diff}  %% MD5 fails from within emacs
\fi

%%%
%%% SVG (Inkscape)
%%%

\ifxetex % chktex 1
\newcommand{\executeiffilenewer}[3]{%
  {\immediate\write18{#3}} % hack
}
\else
\newcommand{\executeiffilenewer}[3]{%
  \ifnum\pdfstrcmp{\pdffilemoddate{#1}}%
    {\pdffilemoddate{#2}}>0%
    {\immediate\write18{#3}}
  \fi%
}
\fi

\newcommand{\includesvg}[2][1.0\textwidth]{%
 \executeiffilenewer{#1.svg}{#1.pdf}%
 {inkscape -z -D --file=#2.svg --export-pdf=#2.pdf --export-latex --export-area-page}%
 \def\svgwidth{#1} 
 \input{#2.pdf_tex}%
} 

%%%
%%% Metropolis Theme
%%%

\usetheme{metropolis}
\metroset{color/block=fill}
\metroset{numbering=none}
\metroset{outer/progressbar=foot}
\metroset{titleformat=smallcaps}

\setbeamercolor{description item}{fg=mLightBrown}
% \setbeamerfont{alerted text}{series=\bfseries}
\setbeamerfont{footnote}{size=\scriptsize}
\setbeamercolor{example text}{fg=mDarkBrown}
\setbeamercolor{block title alerted}{fg=white, bg=mDarkBrown}

\renewcommand*{\UrlFont}{\ttfamily\smaller\relax}

%%%
%%% UTF-8 & Fonts
%%% 

\RequirePackage{unicodesymbols} % after metropolis which loads fontspec

\setmonofont[BoldFont={Cousine Bold},
             ItalicFont={Cousine Italic},
             BoldItalicFont={Cousine Bold Italic}, % Scale=MatchLowercase
             ]{Cousine}

%%%
%%% BibLaTeX
%%%

\RequirePackage[backend=bibtex,
            style=alphabetic,
            maxnames=4,
            citestyle=alphabetic]{biblatex}

\bibliography{local.bib,abbrev3.bib,crypto_crossref.bib,rfc.bib,jacm.bib}

\DeclareFieldFormat{title}{\alert{#1}}
\DeclareFieldFormat[book]{title}{\alert{#1}}
\DeclareFieldFormat[thesis]{title}{\alert{#1}}
\DeclareFieldFormat[inproceedings]{title}{\alert{#1}}
\DeclareFieldFormat[incollection]{title}{\alert{#1}}
\DeclareFieldFormat[article]{title}{\alert{#1}}
\DeclareFieldFormat[misc]{title}{\alert{#1}}

%%% 
%%% Microtype
%%%

\IfFileExists{upquote.sty}{\RequirePackage{upquote}}{}
\IfFileExists{microtype.sty}{\RequirePackage{microtype}}{}

\setlength{\parindent}{0pt}                   %%
\setlength{\parskip}{6pt plus 2pt minus 1pt}  %%
\setlength{\emergencystretch}{3em}            %% prevent overfull lines
\setcounter{secnumdepth}{0}                   %%

%%% Local Variables:
%%% mode: latex
%%% End:
\usepackage{graphicx}
\usepackage{grffile}
\usepackage{longtable}
\usepackage{wrapfig}
\usepackage{rotating}
\usepackage[normalem]{ulem}
\usepackage{amsmath}
\usepackage{textcomp}
\usepackage{amssymb}
\usepackage{capt-of}
\usepackage{hyperref}
\usepackage{microtype}
\usepackage{newunicodechar}
\usepackage[notions,operators,sets,keys,ff,adversary,primitives,complexity,asymptotics,lambda,landau,advantage]{cryptocode}
\usepackage{xspace}
\usepackage{units}
\usepackage{nicefrac}
\usepackage{gensymb}
\usepackage{amsthm}
\usepackage{amsmath}
\usepackage{amssymb}
\usepackage{xcolor}
\usepackage{listings}
\usepackage[color=yellow!40]{todonotes}
\usepackage{filecontents}
\usepackage{url}
\usefonttheme[onlymath]{serif}
\renewcommand{\vec}[1]{\ensuremath{\mathbf{#1}}\xspace}
\newcommand{\bigsoftO}[1]{\ensuremath{\tilde{\mathcal{O}}}\left({#1}\right)\xspace}
\setbeamercolor{example text}{fg=mDarkBrown}
\usetheme{default}
\author{Martin R. Albrecht}
\date{28 March 2018}
\title{The Approximate GCD Problem}
\subtitle{A post-quantum problem that is easier to understand than RSA}
\hypersetup{
pdfauthor={Martin R. Albrecht},
pdftitle={The Approximate GCD Problem},
pdfkeywords={},
pdfsubject={},
pdfcreator={Emacs 25.2.2 (Org mode 9.1.8)},
pdflang={English},
colorlinks,
citecolor=gray,
filecolor=gray,
linkcolor=gray,
urlcolor=gray
}
\begin{document}

\maketitle
\begin{frame}{Outline}
\tableofcontents
\end{frame}


\section{Greatest Common Divisors}
\label{sec:org1749bee}

\begin{frame}[fragile,label={sec:org0c20302}]{Euclidean algorithm}
 Given two integers \(a, b < N = 2^κ\) the Euclidean algorithm computes their greatest common divisor \(\gcd(a,b)\).

\lstset{language=Python,label= ,caption= ,captionpos=b,numbers=none}
\begin{lstlisting}
def gcd(a, b):
    if b == 0:
        return a
    else:
        return gcd(b, a % b)
\end{lstlisting}

The Euclidean algorithm runs in time \(\bigO{κ^2}\).

Best known algorithm runs in time \(\bigO{κ \log^2 κ \log\log κ}\). \footfullcite{ANTS:SteZim04}

For comparison, integer multiplication costs \(\bigO{κ \log κ \log\log κ}\) using the Schönhage–Strassen algorithm.
\end{frame}

\section{RSA}
\label{sec:org73b4419}

\begin{frame}[label={sec:org4cfd7c3}]{Public key encryption}
\begin{description}
\item[{KeyGen}] Bob generates a key pair \((sk, pk)\) and publishes \(pk\).

\item[{Enc}] Alice uses \(pk\) to encrypt message \(m\) for Bob as \(c\).

\item[{Dec}] Bob uses \(sk\) to decrypt \(c\) to recover \(m\).
\end{description}
\end{frame}

\begin{frame}[label={sec:orgbb1aa2e}]{Naive RSA}
\begin{description}
\item[{KeyGen}] The public key is \((N,e)\) and the private key is \(d\), with

\begin{itemize}
\item \(N = p \cdot q\) where \(p\) and \(q\) prime,
\item \(e\) coprime to \(\phi(N) = (p-1)(q-1)\) and
\item \(d\) such that \(e \cdot d \equiv 1 \mod{\phi(N)}\).
\end{itemize}

\item[{Enc}] \(c \equiv m^e \bmod{N}\)

\item[{Dec}] \(m \equiv c^d \equiv m^{e \cdot d} \equiv m^{1} \bmod{N}\)
\end{description}

\begin{block}{Caution}
This naive version of RSA only achieves a very weak form of security — OW-CPA — even against classical adversaries: it is hard to recover random messages.
\end{block}
\end{frame}

\begin{frame}[label={sec:orgd9a7b9b}]{Classical attacks on RSA}
\begin{itemize}
\item An adversary who can factor large integers can break RSA.

\item The best known classical algorithm for factoring is the Number Field Sieve (NFS)

\item It has a \alert{super-polynomial} but \alert{sub-exponential} (in \(\log N\) ) complexity of \[\bigO{e^{1.9 (\log^{1/3} N) (\log\log^{2/3}N)}}\] operations.
\end{itemize}


\pause

\begin{block}{Caution}
This does not mean an adversary \textbf{has} to factor to solve RSA.
\end{block}
\end{frame}

\begin{frame}[label={sec:org7a7fefc}]{Shared factors}
What if two users generate moduli \(N_0 = q_0 ⋅ p\) and \(N_1 = q_1 \cdot p\), i.e. moduli with shared factors?

\begin{itemize}
\item We assume that factoring each of \(N_0\) or \(N_1\) is hard.
\item On the other hand, computing \(\gcd(N_0, N_1)\) reveals \(p\) but costs only \(\bigO{\kappa \log^2 \kappa \log \log \kappa}\) operations when \(N_i \approx 2^\kappa\).
\end{itemize}
\end{frame}

\begin{frame}[label={sec:orgba2f8f7}]{Quantum attacks on RSA}
An adversary with access to a quantum computer with \[ \bigO{\log^2(N) \log\log (N) \log\log\log (N)}\] gates can factor \(N\) using Shor’s algorithm.\footfullcite{FOCS:Shor94}
\end{frame}

\section{The Approximate GCD problem}
\label{sec:org27fab4b}

\begin{frame}[label={sec:orgfaf6116}]{The Approximate GCD problem}
The \alert{Approximate GCD} problem is the problem of distinguishing \[x_i = q_i ⋅ p  \alert{+ r_i}\] from uniform \(\ZZ ∩ [0, X)\) with \(x_i < X\) (\(q_i\), \(r_i\) and \(p\) are secret).
\end{frame}

\begin{frame}[label={sec:org6802ef9}]{The Approximate GCD problem}
\[x_i = q_i ⋅ p  + r_i\]

If \(λ\) is our security parameter (think \(λ=128\)), then

\begin{center}
\begin{tabular}{rrll}
name & sizeof & DGHV10 \footfullcite{EC:DGHV10} & CheSte15 \footfullcite{EC:CheSte15}\\
\hline
\(γ\) & \(x_i\) & \(λ^5\) & \(λ \log λ\)\\
\(η\) & \(p\) & \(λ^2\) & \(λ + \log λ\)\\
\(ρ\) & \(r_i\) & \(λ\) & \(λ\)\\
\end{tabular}

\end{center}
\end{frame}

\begin{frame}[label={sec:org0e238e4}]{Naive encryption}
\begin{description}
\item[{KeyGen}] The public key is \(\{x_i = q_i ⋅ p + 2\,r_i\}_{0 ≤ i < t}\) and the private key is \(p\).

\item[{Enc}] For \(m \in \{0,1\}\) output \(c = m + \sum b_i ⋅ x_i\) with \(b_i \sample \{0,1\}\).

\item[{Dec}] \(m = (c \bmod p) \bmod 2\).
\end{description}

\pause

\begin{block}{Note}
This encryption scheme is not IND-CCA secure but it is IND-CPA secure if the AGCD problem is hard.
\end{block}
\end{frame}

\section{Attacks on the Approximate GCD problem}
\label{sec:org13dc86f}

\begin{frame}[label={sec:org622d0b7}]{Exhaustive search}
Given \(x_0 = q_0 ⋅ p + r_0\) and \(x_1 = q_1 ⋅ p + r_1\) we know that \[p \mid \gcd\left((x_0 - r_0), (x_1 - r_1)\right)\]


Guess \(r_0\) and \(r_1\)!

\begin{block}{Cost}
\(2^{2ρ}\) GCDs
\end{block}
\end{frame}

\begin{frame}[label={sec:org01476e9}]{Exhaustive search + multiplication}
Compute \[\gcd\left(x_0', \prod_{i=0}^{2^ρ-1} (x_1 - i) \bmod x_0'\right)\] for all \(x_0' = x_0 - j\) with \(0 \leq j < 2^{ρ-1}\).

\begin{block}{Cost}
\(2^ρ\) GCDs, \(2^{2ρ}\) multiplications
\end{block}
\end{frame}

\begin{frame}[label={sec:org0fba55b}]{Time-Memory Trade Off}
\begin{lemma}
Assume that we have \(τ\) samples \(x_0 , … ,x_{τ-1}\) of a given prime \(p\), of the hidden form \(x_i = q_i ⋅ p + r_i\), then \(p\) can then be recovered with overwhelming probability in time \(\tilde{\mathcal{O}}(2^{\frac{τ+1}{τ-1}ρ})\). \footfullcite{EC:CorNacTib12}
\end{lemma}
\end{frame}

\begin{frame}[label={sec:org4c10cd1}]{Lattice attacks}
Given \(x_0  = q_0 p + r_0\) and \(x_1  = q_1 p + r_1\), consider

\begin{eqnarray*}
q_0 x_1 - q_1 x_0 & = & q_0 (q_1 p + r_1) - q_1 (q_0 p - r_0)\\
                  & = & q_0 q_1 p + q_0 r_1 - q_1 q_0 p - q_1 r_0\\
& = & q_0 r_1 - q_1 r_0
\end{eqnarray*}

and note that \[q_0 x_1 - q_1 x_0 \ll x_i\]

\pause

\begin{block}{Non-starter?}
We don’t know \(q_i\)!
\end{block}
\end{frame}
\begin{frame}[label={sec:orgf7f462d}]{Lattice attacks}
Consider the matrix 

\[\vec{B} = \begin{pmatrix}
2^{\rho + 1}  & x_1  & x_2   & \cdots  & x_t\\
              & -x_0 &       &         & \\
              &      &  -x_0 &         & \\
              &      &       &  \ddots & \\
              &      &       &         &  -x_0\\
\end{pmatrix}\]

multiplying on the left by the vector \(\vec{q} = (q_0, q_1, q_2, \cdots, q_t)\) gives
\begin{align*}
\vec{v} &= (q_0, q_1, \cdots, q_t) \cdot \vec{B} \\
        &= (q_0\, 2^{ρ+1}, q_0 x_1 - q_1 x_0, \cdots, q_0 x_t - q_t x_0)\\
        &= (q_0\, \alert{2^{ρ+1}}, q_0 \alert{r_1} - q_1 \alert{r_0}, \cdots, q_0 \alert{r_t} - q_t \alert{r_0})
\end{align*}
which is a vector with small coefficients compared to \(x_i\).
\end{frame}

\begin{frame}[label={sec:org0f46ccb}]{Finding short vectors}
Te set of all integer-linear combinations of the rows of \(\vec{B}\) the \alert{lattice} spanned by (the rows of) \(\vec{B}\).

\begin{description}
\item[{SVP}] finding a \alert{shortest} non-zero vector on \alert{general} lattices is NP-hard.

\item[{Gap-SVP\(_\gamma\)}] Differentiating between instances of SVP in which the answer is at most 1 or larger than \(\gamma\) on \alert{general} lattices is a well-known and presumed quantum-hard problem for \(\gamma\) polynomial in lattice dimension.
\end{description}

\begin{block}{Easy SVP}
GCD is SVP on \(\ZZ^2\). For example, \(\vec{B} = {[21, 14]}^T\), \(\vec{v} = (-1,1)\), \(\vec{v} ⋅\vec{B} = 7\).
\end{block}
\end{frame}

\begin{frame}[label={sec:org99f3910}]{Reduction to presumed hard lattice problem}
We can show that an adversary \alert{has} to solve Gap-SVP.

\begin{block}{AGCD → LWE}
If there is an algorithm efficiently solving the AGCD problem then there exists an algorithm which solves the \textbf{Learning with Errors} (LWE) problem with essentially the same performance. \footfullcite{EC:CheSte15} 
\end{block}

\begin{block}{LWE → Gap-SVP}
If there is an algorithm efficiently solving the LWE problem then there exists a quantum algorithm which solves worst-case Gap-SVP instances.\footfullcite{STOC:Regev05}
\end{block}
\end{frame}

\begin{frame}[label={sec:org3f1dbc5}]{Learning with Errors (in normal form)}
Given \((\vec{A},\vec{c})\) with \(\vec{c} \in \ZZ_q^{m}\), \(\vec{A} \in \ZZ_q^{m × n}\), small \(\vec{s} \in \ZZ^{n}\) and small \(\vec{e} \in \ZZ^{m}\) is

\[
\left(\begin{array}{c}
\\
\\
\\ 
\vec{c} \\
\\
\\
\\
\end{array} \right) = \left(
\begin{array}{ccc}
\leftarrow & n & \rightarrow \\
\\
\\ 
& \vec{A} & \\
\\
\\
\\
\end{array} \right) \times \left( \begin{array}{c}
\\
\vec{s} \\
\\
\end{array} \right) + \left(
\begin{array}{c}
\\
\\
\\ 
\vec{e} \\
\\
\\
\\
\end{array} 
\right)
\]

or \(\vec{c} \sample \mathcal{U}(\ZZ_q^{m})\).
\end{frame}

\begin{frame}[label={sec:org11219d7}]{From vectors to scalars}
LWE with modulus \(q^n\) and dimension \(1\) is as hard as LWE with modulus \(q\) and dimension \(1\).

\[q^{d-1} \cdot \Angle{\vec{a},\vec{s}} \approx \left(\sum_{i=0}^{n-1} q^{i} \cdot a_{i}\right) \cdot \left(\sum_{i=0}^{d-1} q^{d-i-1} \cdot s_{i}\right) \bmod q^{d} = \tilde{a} \cdot \tilde{s} \bmod q^{d}.\] 

\begin{block}{Example}
\begin{align*}
\left(a_0 + q\cdot a_1\right) \cdot \left(q \cdot s_0 + \cdot s_1\right) &= q(a_0\cdot s_0 + a_1 \cdot s_1) + (a_1 \cdot s_1) + q^2 (a_1 \cdot s_0)\\
&\equiv q(a_0\cdot s_0 + a_1 \cdot s_1) + (a_1 \cdot s_1) \bmod q^2\\
&\approx q(a_0\cdot s_0 + a_1 \cdot s_1) \bmod q^2
\end{align*}
\end{block}
\end{frame}



\begin{frame}[standout,label={sec:orgff646ac}]{Fin}
\begin{center}
\begin{center}
\includegraphics[width=.9\linewidth]{./kitten-01.jpg}
\end{center}

\alert{\Large Questions?}
\end{center}
\end{frame}

\section{Bonus}
\label{sec:org57bb053}

\begin{frame}[label={sec:orgc88c16d}]{Homomorphic encryption}
Given \(c_i = q_i ⋅ p + m_i'\) with \(m_i' = 2\,r_i + m_i\).
\begin{itemize}
\item We can compute \[c' = c_0 ⋅ c_1 = q_0 q_1 p^2 + q_0 m_1' p  + q_1 m_0' p + m_0' ⋅ m_1'\] to get \(c' \bmod p =  m_0' ⋅ m_1'\) and \(m_0' ⋅ m_1' \bmod 2 = m_0 ⋅ m_1\).
\item We can also compute \[c' = c_0 + c_1 = (q_0 + q_1) p + (m_0' + m_1')\] to get \(c' \bmod p \bmod 2 = m_0 \oplus m_1\).
\end{itemize}

We can compute with encrypted data.\footnote{\url{https://crypto.stanford.edu/craig/easy-fhe.pdf}}
\end{frame}
\end{document}